\section{An analysis pipeline with explicit evidence types}
\label{sec:analysis}

\subsection{Values, bounds, and numerical evidence}
An analysis result should contain its resource model, method, lower/upper
direction, precision, feasibility diagnostics, and evidence. A floating-point
solver status of \texttt{optimal} is not a mathematical certificate. In
particular, the primal objective of a numerical maximization is not, by
itself, a rigorous upper bound on the exact problem.

\paragraph{Classical value.}
Deterministic enumeration is exact for rational verifier tables. The number
of deterministic strategy tuples is
$D=\prod_i |A_i|^{|X_i|}$. Enumerating all but one party and computing that
party's optimal answer separately for each question saves a factor without
changing the optimum. This is still exponential in general. Store a maximizing
response table and its exact rational score. A proof of optimality also needs
the exhaustive-search argument or a separately checked dual/branch certificate.

\paragraph{Non-signalling value.}
Maximize $\langle c,p\rangle$ over positivity, normalization, and
Eq.~\eqref{eq:ns}. This finite linear program gives an upper bound for the
models in Eq.~\eqref{eq:chain}. Marginal constraints must include unsupported
question tuples unless their removal has been justified. Numerical feasibility
residuals are informative but do not replace exact dual feasibility. Rational
reconstruction provides a separate certificate path: reconstruct primal and
dual candidates, verify all equalities and inequalities over the rationals,
and compare their objective values. The implementation can certify a rational
optimum when these bounds coincide, or return the certified upper bound
together with the verified interval when they differ. Both primal and dual
feasibility checks must pass. Failed reconstruction is not reported as a proof.

\paragraph{Quantum outer bounds.}
The NPA hierarchy introduces moments
$\Gamma_{u,v}=\langle u^\dagger v\rangle$ of words in measurement operators,
requires $\Gamma\succeq0$, and imposes the algebraic identities of projectors
and interparty commutation. Finite levels are outer relaxations. The full
hierarchy converges to the commuting-operator set, not in general to
$\Q_{\rm a}$ \cite{npa2008}. Thus the same finite-level bound is valid for
finite-dimensional quantum strategies, but its interpretation must state the
limiting model. Restricted words such as $1+A+B$ and $1+A+B+AB$ are distinct
relaxations, not arbitrary-dimension quantum optimizers.
The implemented level-$1$ relaxation also imposes explicit nonnegative
probabilities, so its precise meaning is $Q_1\cap\NS$; this convention matters
when comparing published bounds.

\paragraph{XOR specialization.}
For a two-party binary-output XOR game with winning rule
$a\oplus b=f(x,y)$, define $C_{xy}=\mu(x,y)(-1)^{f(x,y)}$. The quantum winning
probability is $(1+\beta_q)/2$, where
\begin{equation}
 \beta_q=\max_{\|u_x\|=\|v_y\|=1}\sum_{x,y}C_{xy}u_x\cdot v_y.
 \label{eq:xor}
\end{equation}
The Gram formulation is a semidefinite program, with a quantum realization
available from the vector representation \cite{cleve2004}. This is an exact
characterization of the unrestricted two-party XOR value; a numerical solve
remains approximate. It does not extend unchanged to arbitrary multipartite
XOR tests or arbitrary non-XOR predicates.

\paragraph{Quantum witnesses.}
An explicit density matrix and local POVMs define a feasible strategy by
Eq.~\eqref{eq:quantum}. Check dimensions, Hermiticity, positivity, unit trace,
and POVM completeness before evaluating it. Floating-point checks yield a
numerically validated witness. Symbolic exact checks or interval error bounds
are needed for certified feasibility. Alternating optimization can search for
witnesses in a fixed dimension, but its stationary point is a lower-bound
candidate, not a global optimum.
The implemented bipartite see-saw alternates a largest-eigenvector state
update with local POVM semidefinite updates and uses seeded restarts.

\subsection{Symbolic certificates: CHSH as a complete example}
Use observables $A_0,A_1,B_0,B_1$ with $A_x^2=B_y^2=\id$ and
$[A_x,B_y]=0$. Let
$\mathcal B=A_0(B_0+B_1)+A_1(B_0-B_1)$. Direct expansion yields
\begin{equation}
 2\sqrt2\id-\mathcal B
 =\frac1{\sqrt2}\left[
 \left(A_0-\frac{B_0+B_1}{\sqrt2}\right)^2+
 \left(A_1-\frac{B_0-B_1}{\sqrt2}\right)^2\right].
 \label{eq:chsh-sos}
\end{equation}
Each squared Hermitian operator is positive, proving
$\langle\mathcal B\rangle\leq2\sqrt2$. The state
$\ket{\Phi^+}=(\ket{00}+\ket{11})/\sqrt2$, observables $A_0=Z,A_1=X$,
$B_0=(Z+X)/\sqrt2,B_1=(Z-X)/\sqrt2$, attain the bound. For uniform questions
and winning rule $a\oplus b=xy$,
\[
 \val_{\Lset}=\tfrac34,\qquad
 \val_{\Q_{\rm fin}}=\tfrac12+\tfrac{\sqrt2}{4},\qquad
 \val_{\NS}=1.
\]
The classical bound follows because the four required parity equations have
inconsistent XOR and three can be satisfied. The non-signalling value is
attained by $p(a,b\mid x,y)=1/2$ when $a\oplus b=xy$, and zero otherwise.
Equation~\eqref{eq:chsh-sos} proves the quantum upper bound independently of
solver tolerances. SOS methods also expose the operator relations needed for
self-testing, although extraction requires further arguments
\cite{bamps2015}.

\subsection{Metrics beyond one optimal value}
\label{sec:benchmarks}
The following exact literature values fix the conventions of the benchmark
constructors. They are mathematical expectations, not a claim that every
quantum value is certified by the implemented numerical backends.
\begin{center}
\begin{tabular}{lccc}
\toprule
Verifier and sampling & $\val_{\Lset}$ & $\val_{\Q_{\rm fin}}$ & $\val_{\NS}$\\
\midrule
Uniform CHSH & $3/4$ & $(2+\sqrt2)/4$ & $1$\\
Three-party GHZ, uniform even inputs & $3/4$ & $1$ & $1$\\
Magic square, uniform row/column & $8/9$ & $1$ & $1$\\
Three-party GYNI, uniform even inputs & $1/4$ & $1/4$ & $1/3$\\
\bottomrule
\end{tabular}
\end{center}
GHZ requires $\sum_i a_i=\frac12\sum_i x_i\pmod2$ under the even-input
promise; the three-party case of the parity-game analysis gives its values
\cite{brassard2003}. Magic square uses an even row for Alice, an odd column
for Bob, and agreement at their intersection \cite[Sec.~5]{brassard2005}.
This is the row/column version, not the row-or-column/single-cell variant
whose classical value is $17/18$ \cite[Sec.~3.3]{cleve2004}.
The perfect quantum strategies imply the stated non-signalling values.
GYNI accepts iff $(a_1,a_2,a_3)=(x_2,x_3,x_1)$. Almeida et al. prove that
quantum resources give no advantage for any input prior, and derive the
non-signalling value $1/3$ for this specific even-input prior
\cite[Eqs.~(10)--(12)]{almeida2010}. GYNI therefore tests an important failure
mode: a non-signalling gap alone does not demonstrate a quantum advantage.

For normalized games, report the additive gap
$\Delta=\val_Q-\val_L$ and the actual score conventions. Ratios are unstable
under harmless offsets and should not be compared without a normalization.
For a specified noise model
$p_v=vp_*+(1-v)p_0$ with $S_G(p_*)>\val_L>S_G(p_0)$, the visibility threshold
for violating the local bound is
\[
 v_{\rm crit}=\frac{\val_L-S_G(p_0)}{S_G(p_*)-S_G(p_0)}.
\]
This is a threshold for that family and witness, not for all noisy states or
all possible measurements. Loss, biased inputs, misalignment, and adversarial
noise require separate models. Conditioning on successful detection can change
the local bound and cannot be treated as unconditional performance.

Other useful questions include dimension/entanglement witnesses, nonlocal
membership, distance to the local set, robust extraction, randomness against
quantum side information, and sample complexity. These require different
optimization objects and assumptions. A game-value gap alone supplies neither
a device-independent entropy bound nor cryptographic security. Likewise,
finite-sample significance requires a sampling and memory model, explicit
confidence levels, and a valid statistical test; point estimates are not
confidence bounds.

\subsection{Why no complete automatic analyzer can exist}
There are finite games whose entangled value is not attained in finite
dimension \cite{slofstra2019}. More decisively, the gap problem distinguishing
entangled value $1$ from value at most $1/2$ is undecidable in general
\cite{ji2020}. A universal finite game language is therefore compatible with
the impossibility of a uniformly terminating, complete value-analysis
algorithm. The engineering consequence is typed partial procedures: exact
algorithms on restricted fragments, certified bounds when available, numerical
evidence with explicit limits, and an honest unresolved interval otherwise.
